% from ega4-9-fr.tex
\oldpage[IV-3]{68}

$2^\circ$ $\eta\in E$, donc $Z_\eta$ est dense dans $Z'_\eta$. Montrons
d'abord que l'on peut supposer $Z'$ \emph{fermé}. En effet, $Z_\eta$ est dense
dans $\overline{Z'_\eta}$ (adhérence prise dans $X_\eta$)~; posons
$V_\eta=X_\eta-\overline{Z_\eta}$, qui est ouvert dans $X_\eta$ et ne
rencontre pas $Z'_\eta$~; on peut supposer $V_\eta$ de la forme $V\cap X_\eta$,
où $V$ est ouvert (donc constructible) dans $X$, et l'hypothèse
$V_\eta\cap Z'_\eta=\varnothing$ entraîne alors
$V_s\cap Z'_s=\varnothing$ pour tout $s$ voisin de $\eta$ en vertu de (9.5.1).
Remplaçant $S$ par un voisinage ouvert de $\eta$, on peut donc supposer que
$V\cap Z'=\varnothing$, donc $V\cap\overline{Z'}=\varnothing$ (adhérence prise
dans $X$), et par suite $(\overline{Z'})_\eta=\overline{Z'_\eta}$, d'où notre
assertion. L'ensemble $Z'$ est alors réunion de ses composantes irréductibles
en nombre fini, et en restreignant encore $S$ à un voisinage de $\eta$, on
peut supposer que toutes les composantes irréductibles $Z'_i$ de $Z'$
rencontrent $X_\eta$, d'où résulte que $X_\eta$ contient les points génériques
des $Z'_i$ $(\mathbf 0_{\mathrm I},2.1.8)$. Dire que $Z_s$ est dense dans
$Z'_s$ équivaut alors à dire que chacun des $(Z\cap Z'_i)_s$ est dense dans
$(Z'_i)_s$, et l'on est ainsi ramené au cas où $Z'$ est irréductible.
Remplaçant alors au besoin $X$ par le sous-préschéma réduit ayant $Z'$ pour
espace sous-jacent, on voit qu'on peut supposer que $Z'=X$ et que $X$ est
\emph{intègre} et domine $S$. Enfin, en recouvrant $X$ par un nombre fini
d'ouverts affines $W_j$ et remplaçant $Z$ par $Z\cap W_j$, on peut supposer
que $X=\operatorname{Spec}(A)$, où $A$ est un anneau noethérien \emph{intègre}.
Comme $X_\eta$ est noethérien intègre et que $Z_\eta$ est constructible dans
$X_\eta$ et dense dans $X_\eta$, $Z_\eta$ contient un ouvert non vide de
$X_\eta$ $(\mathbf 0_{\mathrm{III}},9.2.2)$, que l'on peut supposer de la forme
$(D(t))_\eta$, où $t$ est un élément $\neq0$ de $A$. En remplaçant au besoin
$S$ par un voisinage de $\eta$, on peut en outre, en vertu de la relation
$(D(t))_\eta\subset Z_\eta$, supposer que $D(t)\subset Z$ (9.5.2). Enfin,
comme l'homothétie de rapport $t_\eta$ dans $\mathcal O_{X_\eta}$ est
injective, il résulte de (9.4.5) que pour $s$ voisin de $\eta$, $t_s$ est
$\mathcal O_{X_s}$-régulier, donc $(X_s)_{t_s}$ est dense dans $X_s$, et
\emph{a fortiori} il en est de même de $Z_s$ qui contient $(X_s)_{t_s}$.

\begin{corollary}[9.5.4]
\label{IV.9.5.4-fr}
Soient $f:X\to S$ un morphisme de présentation finie, $Z$ une partie
localement constructible de $X$. L'ensemble $E$ des $s\in S$ tels que $Z_s$
soit fermé \textnormal{(}resp. ouvert, resp. localement fermé\textnormal{)}
dans $X_s$ est localement constructible dans $S$.
\end{corollary}

Dire que $Z_s$ est ouvert dans $X_s$ signifie que
$(X-Z)_s=X_s-Z_s$ est fermé dans $X_s$, et comme $X-Z$ est localement
constructible, on peut se borner à considérer l'ensemble des $s\in S$ tels que
$Z_s$ soit fermé et l'ensemble des $s\in S$ tels que $Z_s$ soit localement
fermé.

Vérifions encore dans chaque cas les deux conditions de (9.2.2, (ii)). La
première résulte du fait que $p:X_{(k')}\to X$ est fidèlement plat et
quasi-compact, et de (2.3.12) et (2.3.14). Vérifions donc la seconde condition,
$S$ étant supposé affine, noethérien et intègre, de point générique $\eta$.
Posons $Z'=\overline Z$~; le même raisonnement que dans (9.5.3) montre que
$Z'_\eta$ est égal à l'adhérence de $Z_\eta$ dans $X_\eta$~; en vertu de
(9.5.3), il y a donc un voisinage $U$ de $\eta$ tel que pour $s\in U$, $Z_s$
soit dense dans $Z'_s$, ce dernier étant fermé dans $X_s$. Dire que $Z_s$ est
fermé dans $X_s$ signifie alors que $Z''_s=\varnothing$, où $Z''=Z'-Z$~; il
résulte donc de (9.5.1) que l'ensemble $E$ des $s\in S$ où
$Z''_s=\varnothing$ est tel que $E$ ou $S-E$ contienne un voisinage de
$\eta$. Dire que $Z_s$ est localement fermé dans $X_s$ signifie que $Z''_s$
est fermé dans $X_s$~; il suffit donc d'appliquer le résultat précédent en
remplaçant $Z$ par $Z''$, qui est localement constructible dans $X$.

